\section{Extended Related Work}
\label{app:related_work}

% ===========================================================================
% Paradigm 1: Grammar-Constrained Decoding
% ===========================================================================
\paragraph{Grammar-Constrained Decoding.}
Recent work has established efficient algorithms for restricting LLM outputs to conform to context-free grammars, making grammar-constrained decoding the de facto standard for obtaining structured output from language models.
XGrammar~\cite{dong2025xgrammar} partitions the vocabulary into context-independent and context-dependent tokens and maintains persistent pushdown automata across decoding steps, achieving up to 100$\times$ speedup over prior CFG-constrained decoding methods.
Outlines~\cite{willard2023outlines} reformulates constrained generation as finite-state machine transitions with offline vocabulary indexing, which enables efficient enforcement of regular-expression and JSON-schema constraints without runtime overhead per token.
LMQL~\cite{beurer-kellner2023lmql} introduces a declarative query language that compiles type, length, and regex constraints into token-level masks, providing a programmable interface for specifying multi-variable output structures.
These systems ensure that every generated token belongs to a grammar-valid continuation, but they do not address the \emph{distribution distortion} introduced by naive token masking.
A parallel line of work tackles this problem directly: DOMINO~\cite{beurer-kellner2024domino} minimizes distributional divergence during CFG enforcement by formulating mask selection as an optimization problem, and ASAp~\cite{park2024grammar} provides a provably unbiased sampling algorithm that preserves the original LM distribution under grammar constraints.
Together, these advances have made grammar-constrained decoding both fast and statistically principled, and provide a robust foundation for structured generation.
However, all of these methods operate exclusively at the syntactic level: they ensure format validity (e.g., well-formed JSON, valid enum values) but cannot enforce cross-field semantic constraints such as ontology domain/range rules that depend on previously generated content.
This syntactic-only limitation motivates OntoCoDec's semantic extension: our grammar layer builds on XGrammar for structural (JSON schema and enum) constraints, while a second semantic layer enforces relational type constraints that are inherently beyond the expressive power of context-free grammars.

% ===========================================================================
% Paradigm 1.5: Semantic Constrained Decoding for Code
% ===========================================================================
\paragraph{Semantic Constrained Decoding for Code.}
In programming-language generation, constrained decoding has been extended beyond syntax to enforce semantic properties such as type correctness; this work demonstrates that richer-than-CFG constraints can be compiled into token-level filtering.
In the text-to-SQL domain, PICARD~\cite{scholak2021picard} maintains an incremental SQL parsing state that rejects tokens violating syntax or referencing non-existent database columns at each decoding step, achieving state-of-the-art results on Spider and CoSQL without modifying model weights.
While mechanistically similar to OntoCoDec's token-level filtering, PICARD operates over a \emph{fixed, fully enumerated} schema (database table and column names known before decoding begins).
In contrast, OntoCoDec's entity registry is populated \emph{during} generation as entity spans are extracted from free text; constraint sets evolve with the decoding prefix, a fundamental distinction between querying a known schema and extracting from unstructured text.
Synchromesh~\cite{synchromesh_iclr2022} pioneers Constrained Semantic Decoding (CSD), which enforces syntax, type, and scope constraints simultaneously during code generation from pre-trained LMs without fine-tuning; this establishes that semantic constraints can be expressed as dynamic token masks conditioned on the generation prefix.
MGD~\cite{mgd_neurips2023} uses static-analysis monitors to perform incremental type checking at decode time; SantaCoder-1.1B with MGD outperforms text-davinci-003 on repository-level code completion by ensuring that every generated identifier is type-correct with respect to the surrounding code context.
Most recently, \citet{pldi2025_type_constrained_codegen} compile full type systems, including generic types and subtyping, into prefix automata for token-level logits masking and reduce compilation errors by over 50\%.
This last system is the closest mechanistic analog to OntoCoDec, applying the same ``type system $\to$ automaton $\to$ logits mask'' pipeline, though in a fundamentally different domain.
The key technical difference lies in the constraint structure: programming-language type systems feature deep subtyping hierarchies and parametric polymorphism, whereas IE ontology axioms specify flat domain/range sets; a relation $R$ accepts head entities from set $\mathcal{D}(R)$ and tail entities from set $\mathcal{R}(R)$, with no subtype chains.
Conversely, IE introduces a challenge absent in code generation: the constraint target (entity type) must be inferred from free-text spans rather than read from declared type annotations, so a stateful entity registry must track extracted types during autoregressive generation.
These methods collectively demonstrate that semantic (type) constraints at decode time are both feasible and highly effective, but exclusively in the code domain where type systems are formally specified.
OntoCoDec is the first to \emph{automatically compile} IE ontology axioms, including domain/range constraints, into token-level decoding masks for information extraction, migrating the type-constrained decoding paradigm from code generation to IE.

% ===========================================================================
% Paradigm 2: Vocabulary/Type-Constrained Generation
% ===========================================================================
\paragraph{Vocabulary and Type-Constrained Generation.}
The idea of enforcing entity-type compatibility constraints for information extraction dates back to pre-neural methods.
\citet{roth_yih_2004} formulate joint entity and relation identification as an integer linear program (ILP) that enforces type constraints (e.g., \textsc{Emp-Org} only between \textsc{Per} and \textsc{Org} entities); this establishes the foundational principle that relational type compatibility should be a hard constraint, not a soft preference learned from data.
In the autoregressive era, GENRE~\cite{decao2021genre} pioneers trie-based constrained beam search for autoregressive entity retrieval, restricting each decoding step to tokens consistent with valid entity names from a predefined catalog.
GenIE~\cite{josifoski2022genie} extends this approach to end-to-end information extraction with state-based constraints that enforce valid triple structures, achieving strong performance on closed-world knowledge base population, though it requires task-specific fine-tuning on the target schema.
AdaKGC~\cite{adakgc_emnlp2023} introduces schema-constrained dynamic decoding with trie-based prefix trees that adapt to the target schema, enforcing that generated entity and relation names belong to a predefined inventory, a vocabulary-level constraint that ensures name validity but does not verify type compatibility between a relation and its arguments.
GCD~\cite{geng2023gcd} provides a unified grammar-constrained decoding framework for structured NLP tasks including IE, introducing input-dependent grammars that adapt constraints based on the source text. GCD is the closest precursor to OntoCoDec in applying logits masking to IE, though limited to syntactic constraints.
NeuroLogic Decoding~\cite{lu2021neurologic} expresses lexical inclusion and exclusion constraints as predicate logic formulas enforced during beam search, offering a general mechanism for token-level control that can encode vocabulary membership but not relational type rules.
These approaches constrain \emph{which tokens} can appear (vocabulary membership) but not \emph{which type combinations} are valid for a given relation.
OntoCoDec's semantic layer enforces relational constraints (e.g., that the head entity of relation $R$ must belong to type set $\mathcal{D}(R)$) which cannot be expressed as static vocabulary restrictions because the valid set depends on the relation type generated earlier in the same sequence.
While Roth \& Yih's ILP formulation shares the same type-compatibility motivation, it operates as a post-hoc reranker over discrete candidates; OntoCoDec integrates enforcement directly into autoregressive decoding, preventing violations rather than filtering them.

% ===========================================================================
% Paradigm 3: Structural Constraints for IE
% ===========================================================================
\paragraph{Structural Constraints via Training and Prompting.}
An orthogonal approach encodes IE output structure through model training or prompt design rather than decoding-time enforcement, relying on learned representations to produce schema-compliant outputs.
UIE~\cite{lu2022uie} defines a Structured Extraction Language and trains a unified text-to-structure model across 13 IE datasets, relying on large-scale pre-training to internalize output format constraints.
GoLLIE~\cite{sainz2024gollie} expresses annotation guidelines as Python dataclass definitions and fine-tunes Code-LLaMA to generate typed code, achieving strong zero-shot transfer to unseen IE schemas through the implicit type structure of programming languages.
KnowCoder~\cite{knowcoder_acl2024} takes this further by encoding 29K entity types from Wikidata as a Python class hierarchy, training a Code-LLM to generate structured extraction code that implicitly respects type inheritance, the most comprehensive schema-as-code approach to date.
CodeIE~\cite{li2023codeie} similarly reformulates IE as code generation, using the structural regularity of programming languages to reduce output format errors.
ADELIE~\cite{qi2024adelie} uses supervised fine-tuning (SFT) and Direct Preference Optimization (DPO) to align LLMs on IE tasks, achieving state-of-the-art results among open-source models while preserving general capabilities, though compliance with ontological constraints remains a learned soft preference rather than a hard constraint.
In the supervised setting, joint entity-relation architectures such as SpERT~\cite{eberts2020spert} and PFN~\cite{yan2021pfn} enforce type consistency implicitly through shared representations: SpERT classifies entity spans and relation pairs in a single BERT pass with span-level negative sampling, while PFN decomposes feature encoding into entity-specific, relation-specific, and shared partitions via gating mechanisms.
These models achieve strong performance but are limited to fixed, predefined schemas and cannot generalize to unseen ontologies without retraining, precisely the ontology-portable scenario requiring no constraint-specific training that OntoCoDec targets.
EnriCo~\cite{zaratiana2024enrico} applies Answer Set Programming (ASP) constraints during decoding of joint entity-relation extraction models, enforcing structural type compatibility at generation time. This is a decoding-time approach, though with model-specific, manually authored constraints rather than ontology-driven compilation.
Training-based methods provide \emph{soft} structural guidance: the model learns to prefer valid outputs but cannot enforce compliance at inference time, as demonstrated by non-trivial violation rates even in fine-tuned systems.
Importantly, training-based and decoding-time approaches are complementary rather than competing: a system like GoLLIE could be combined with OntoCoDec's constraint layers, using fine-tuning to maximize the probability of valid outputs while decoding-time enforcement eliminates residual violations.
OntoCoDec provides \emph{hard} decoding-time enforcement that requires no constraint-specific training, is automatically compiled from IE ontologies, and composable with any of these training-based approaches.

% ===========================================================================
% Paradigm 4: Type-Checking / Ontology Validation
% ===========================================================================
\paragraph{Post-hoc Ontology Validation.}
Rather than constraining generation, several methods validate LLM outputs against ontological rules after generation, applying a detect-then-fix paradigm.
ODKE+~\cite{khorshidi2025odke} dynamically generates ontology snippets containing domain/range constraints as prompt context, achieving 98.8\% precision across 195 predicates at web scale, but constraint adherence relies on prompt-level guidance, not enforcement, leaving residual violations unaddressed.
SPIRES/OntoGPT~\cite{spires_2024} provides an ontology-based extraction framework that grounds LLM outputs against OBO ontologies post-generation: it validates extracted entities and relations through ontology lookup and corrects mismatches in a second pass, the prompt-then-validate paradigm in the biomedical domain.
OntoLogX~\cite{ontologx_2025} applies Shapes Constraint Language (SHACL) validation to LLM-generated knowledge graphs, iteratively re-prompting the LLM to correct constraint violations and reducing them by approximately 10$\times$ across multiple correction rounds. This demonstrates the detect-then-fix paradigm's effectiveness but also its inherent multi-pass cost.
\citet{upadhyay2023openie} inject entity-type knowledge as constraints into Open IE via both constrained inference and constraint-aware training, improving F1 by 29\% on CaRB, though their constrained inference operates over a discrete candidate set rather than during autoregressive generation.
\citet{regino2025kgcurator} propose a multi-stage LLM pipeline for validating Resource Description Framework (RDF) triple insertions against RDF Schema (RDFS) domain/range constraints, where each validation step requires a separate LLM call, effective but computationally expensive at scale.
RELATE~\cite{olasunkanmi2025relate} addresses the alignment gap between LLM-extracted biomedical relations and standardized ontology predicates through a three-stage pipeline: ontology preprocessing with predicate embeddings, similarity-based retrieval via SapBERT, and LLM-based reranking with negation handling, achieving 94\% accuracy@10 on ChemProt.
While RELATE demonstrates effective ontology grounding for relation extraction, its constraints are applied \emph{post-extraction} through retrieval and reranking rather than during generation, and it requires domain-specific embedding models trained on biomedical corpora.
OntoCoDec enforces ontology constraints \emph{during} autoregressive decoding and requires no additional trained components beyond the base LLM.
Concurrently, \citet{zhou2026ontologytools} compile Web Ontology Language (OWL) ontologies into executable tool interfaces within The World Avatar framework; LLM agents must invoke ontology-aware tools that enforce semantic constraints during knowledge graph construction from scientific text.
This ``ontology-to-tools'' approach shares OntoCoDec's philosophy of compiling formal ontological knowledge into runtime enforcement mechanisms, but operates at the tool-call level (constraining which API functions the agent invokes) rather than at the token level (constraining which tokens the decoder produces), and requires an agentic multi-step workflow rather than single-pass autoregressive generation.
\citet{ekaw2024_ontology_beam} introduce ontology-constrained beam scoring for clinical summarization, applying soft biases (not hard masks) during beam search to steer outputs toward ontology-compliant terms, a middle ground between unconstrained generation and hard enforcement that reduces but cannot eliminate violations.
Post-hoc validation detects violations after generation and requires rejection or correction cycles that waste compute. In our setting, 15.37\% of generated triples contain domain/range violations under unconstrained decoding, a substantial waste of generation effort.
OntoCoDec \emph{prevents} violations during decoding, redirecting probability mass to valid alternatives rather than discarding invalid outputs, with only marginal overhead from logits masking compared to the cost of regenerating or correcting violated triples.
Soft-constraint approaches such as ontology-weighted beam scoring reduce violation frequency but cannot eliminate it, whereas OntoCoDec substantially reduces violations (DR$\downarrow$94\% on SciERC, $\downarrow$45--90\% on DocRED depending on model).

% ===========================================================================
% Paradigm 5: Constrained Decoding Quality Concerns
% ===========================================================================
\paragraph{Quality Impact of Constrained Decoding.}
A growing body of work raises concerns that constrained decoding degrades LLM generation quality, particularly for reasoning-intensive tasks where format enforcement interferes with chain-of-thought processes.
\citet{tam2024letmespeak} conduct a comprehensive evaluation showing that rigid format constraints (JSON, XML) harm performance on math and multi-hop QA benchmarks but \emph{benefit} classification-style tasks such as slot filling, a distinction directly relevant to IE, which is fundamentally an extraction (classification-like) task rather than a reasoning task.
CRANE~\cite{banerjee2025crane} provides theoretical analysis of why restricting outputs to narrow grammars suppresses chain-of-thought reasoning by eliminating intermediate reasoning tokens, and proposes an alternating strategy between unconstrained reasoning phases and constrained output phases to mitigate this effect.
\citet{schall2025hiddencost} conduct a systematic study across 11 models and multiple benchmarks, finding 10--30\% semantic correctness degradation on reasoning tasks under structural constraints, with larger models showing greater absolute degradation.
\citet{zhou2026snowballing} identify ``structure snowballing,'' a failure mode where format compliance cognitive load traps models in repetitive structural patterns rather than improving reasoning during self-correction iterations.
These findings establish a legitimate concern for constrained decoding in reasoning-heavy applications.
However, information extraction occupies a fundamentally different regime: the model extracts entities and relations from input text rather than generating novel reasoning chains, and \citet{tam2024letmespeak}'s own results show that extraction-style tasks are robust to, or even benefit from, format constraints.
OntoCoDec addresses this concern empirically: domain/range violations decrease by 94\% on SciERC (15.37\%$\to$0.93\%) and 45--90\% on DocRED across two models, with F1-neutral impact on SciERC ($+$0.13pp, $p{=}0.58$) and significant R-F1 gains on DocRED (Qwen $p{<}0.001$, Llama $p{=}0.013$).
Critically, OntoCoDec's constraints target only type tokens (entity types, relation labels), leaving free-text spans (entity surface forms and their boundaries) entirely unconstrained, which minimizes the distribution distortion that causes quality degradation in reasoning tasks.
A two-pass variant (S2), explored in ablation (Section~5), separates extraction and constraint enforcement into distinct passes, analogous to CRANE's alternating strategy; however, our primary single-pass approach already avoids reasoning interference by constraining only type tokens while leaving free-text spans unconstrained.

% ===========================================================================
% Our Approach
% ===========================================================================
\paragraph{Our Approach.}
OntoCoDec occupies a unique position in this landscape: it is the first system to \emph{automatically compile} IE ontologies, including domain/range axioms, into token-level semantic masks, extending grammar-constrained decoding from syntactic to semantic enforcement without manually authored rules.
Unlike vocabulary-constrained methods that restrict individual tokens to a fixed set, OntoCoDec enforces \emph{cross-field relational constraints} through a stateful entity registry that tracks extracted types during autoregressive generation; constraint evaluation is dynamic and depends on previously decoded content.
This approach is orthogonal to supervised IE methods and applicable to autoregressive LMs with logits access without model modification or additional training (evaluated on 7B--70B models across three tokenizer families).
Unlike training-based IE systems, OntoCoDec requires no constraint-specific training and substantially reduces type violations at inference time (DR$\downarrow$94\% on SciERC, $\downarrow$45--90\% on DocRED); unlike post-hoc validation, it prevents violations rather than detecting them, redirecting probability mass to valid alternatives with significant R-F1 gains on DocRED ($p \leq 0.013$) and F1-neutral impact on SciERC.
The two constraint layers, grammar (structural) and semantic (relational), compose via bitwise AND with a scale-dependent interaction: the grammar layer ensures every generated type token belongs to the valid vocabulary, which enables the semantic layer to look up and enforce domain/range rules for that type.
